% Chapter 0: 引言 - 文献概述
% Schubert 多项式与 Positivity

\chapter*{引言：文献概述}
\addcontentsline{toc}{chapter}{引言：文献概述}
\label{chap:Intro}

\section{总述}

本笔记学习的主题是 \textbf{Schubert 多项式的 Positivity（正性）问题}。这是代数组合与代数几何中的核心问题之一，与旗流形的上同调理论密切相关。

Schubert 多项式是旗流形 $Fl_n(\mathbb{C})$ 上 Schubert 类的多项式代表元。它们由 Lascoux 和 Schützenberger 引入，具有丰富的代数和组合结构。

本笔记涉及四篇关于 Schubert 多项式正性问题的研究论文，从经典情形逐步推广到量子情形。此外，还有若干参考书籍和补充材料作为基础知识来源。

\section{研究目标}
\label{sec:ResearchGoal}

\textbf{研究课题}：三重 Schubert 演算的 Graham Positivity 猜想

\textbf{核心问题}：证明三重（double）Schubert 多项式乘积展开系数的正性

\textbf{文献组合策略}：根据熊瑞的观察，Chapter 1（Gao-Xiong 2025）和 Chapter 4（Mihalcea 2007）这两篇论文组合在一起，基本上能够搞定三重 Schubert 演算的正性猜想。

\section{术语对照表}
\label{sec:TerminologyTable}

为统一全笔记的术语使用，本节建立以下术语对照表：

\begin{center}
\begin{tabular}{p{4cm} p{6cm} p{4cm}}
\toprule
术语 & 含义 & 参考文献 \\ \midrule
\textbf{Graham Positivity} & 原始情形，要求 $u, v$ 都是 Grassmannian 排列，系数 $c^w_{u,v} \in \mathbb{Z}_{\geq 0}[t_i-y_j]$ & \cite{Gr} \\
\textbf{Refined Graham Positivity} & 推广到满足 separated descents 条件的任意排列 & \cite{GX2025} \\
\textbf{Triple Schubert Positivity} & Gao-Xiong 2025 的主要结果，处理三组变量 $(\mathbf{x}; \mathbf{y}; \mathbf{t})$ 情形 & \cite{GX2025} \\
\textbf{Separated Descents} & 条件 $\max\operatorname{Des}(u) \leq \min\operatorname{Des}(v)$ & \cite{Sa} \\
\bottomrule
\end{tabular}
\end{center}

\textbf{变量约定}：
\begin{itemize}
\item $\mathbf{x} = (x_1, x_2, \ldots)$：旗流形 $Fl_n$ 的坐标变量
\item $\mathbf{y} = (y_1, y_2, \ldots)$：来源 Schubert 类的等变参数
\item $\mathbf{t} = (t_1, t_2, \ldots)$：目标 Schubert 类的等变参数
\end{itemize}

\section{文献摘要总结}
\label{sec:LiteratureSummary}

\subsection{Gao-Xiong 2025: Triple Schubert Positivity}
\textbf{作者}: Yibo Gao, Rui Xiong \\
\textbf{arXiv}: \url{https://arxiv.org/abs/2506.09421} \\
\textbf{年份}: 2025 \\
\textbf{篇幅}: 6 页

\textbf{核心贡献}:
\begin{itemize}
\item 证明了 Samuel 关于三重变量集的 Graham positivity 猜想
\item 建立了 refined 版本的 Graham positivity 定理
\item 作为推论，证明了 Kirillov 关于 skew 除差算子的正性猜想
\end{itemize}

\textbf{摘要}: 本文证明了 Samuel 关于三个变量集中两个双重 Schubert 多项式展开系数的 Graham positivity 猜想。通过建立 Graham positivity 定理的 refined 版本，作为推论证明了 Kirillov 关于 skew 除差算子作用于 Schubert 多项式的正性猜想。

\textbf{关键词}: Graham positivity, double Schubert polynomials, triple Schubert calculus, divided difference operators

\subsection{Samuel 2024: Molev-Sagan Formula}
\textbf{作者}: Matthew J. Samuel \\
\textbf{arXiv}: \url{https://arxiv.org/abs/2401.11060} \\
\textbf{年份}: 2024 \\
\textbf{篇幅}: 30 页

\textbf{核心贡献}:
\begin{itemize}
\item 给出了计算双重 Schubert 多项式乘积的 Molev-Sagan 类型公式
\item 提供了 Pieri 公式
\item 证明了倾斜 Schubert 多项式系数的非负性
\end{itemize}

\textbf{摘要}: 本文给出了计算两组不同系数变量的两个双重 Schubert 多项式乘积的 Molev-Sagan 类型公式。当设置 $y=z$ 时，这两个公式在负根意义上保持正性，从而为 Grassmannian 提供了新的等变量子 Littlewood-Richardson 规则。

\textbf{关键词}: Molev-Sagan formula, double Schubert polynomials, Pieri formula, Littlewood-Richardson rule

\subsection{Kirillov-Maeno 1996: Quantum Schubert Polynomials}
\textbf{作者}: Anatol N. Kirillov, Toshiaki Maeno \\
\textbf{arXiv}: \url{https://arxiv.org/abs/q-alg/9610022} \\
\textbf{年份}: 1996 \\
\textbf{篇幅}: 52 页

\textbf{核心贡献}:
\begin{itemize}
\item 引入了量子双重 Schubert 多项式 $\tilde{S}_w(x, y)$
\item 研究了旗流形的等变量量子上同调
\item 基于量子 Cauchy 恒等式的方法
\item 给出了 Vafa-Intriligator 公式的高 genus 推广
\end{itemize}

\textbf{摘要}: 本文研究了旗流形的等变量量子上同调代数。引入并研究了量子双重 Schubert 多项式 $\tilde{S}_w(x, y)$，这是旗流形等变量量子上同调类的 Lascoux-Schützenberger 型代表元。

\textbf{关键词}: quantum Schubert polynomials, equivariant quantum cohomology, flag variety, Cauchy identity

\subsection{Mihalcea 2007: Equivariant Quantum Positivity}
\textbf{作者}: Leonardo Constantin Mihalcea \cite{Mihalcea} \\
\textbf{arXiv}: \url{https://arxiv.org/abs/math/0603438} \\
\textbf{年份}: 2007 \\
\textbf{篇幅}: 15 页

\textbf{核心贡献}:
\begin{itemize}
\item 证明了 Peterson 猜想（由 Graham 证明）在等变量量子上同调中仍然成立
\item 推广了经典上同调的正性结果到量子版本
\end{itemize}

\textbf{摘要}: D. Peterson 的猜想（由 W. Graham 证明）指出齐次空间 $G/P$ 的 $T$-等变上同调（equivariant cohomology）的结构常数满足某种正性性质。本文证明了在更一般的等变量量子上同调情况下，这个正性性质仍然成立。

\textbf{关键词}: equivariant quantum cohomology, positivity, Schubert calculus, Gromov-Witten invariants

详见 \cref{chap:EquivariantQuantumPositivity}（Chapter 4）。

\section{补充参考资料}

以下文献作为基础知识的补充来源：

\begin{itemize}
\item \textbf{Anderson-Fulton}: \textit{Equivariant Cohomology in Algebraic Geometry} — 等变量子上同调的权威教材，详见 \cref{def:EquivariantQuantumCohomology}（Chapter 4）
\item \textbf{Li 2023}: \textit{Schubert Calculus -- A Brief Introduction} — Schubert calculus 入门讲义
\item \textbf{Buch 2001}: \textit{Quantum Cohomology of Grassmannians} — Grassmannian 量子上同调环的简洁证明
\item \textbf{Fan-Gao-Gu-Xiong 2025}: \textit{Bumpless Pipe Dreams Meet Puzzles} — 旗流形上同调的全面综述，涵盖 Pipe Dreams 和 Bumpless Pipe Dreams 的组合理论（详见 Chapter 2）
\end{itemize}

\section{学习路径与前置基础知识}

根据上述文献的关系，建议学习顺序如下：

\begin{enumerate}
\item \textbf{首先}: 学习 Gao \& Xiong (2025) — 理解 triple positivity 的证明
\item \textbf{其次}: 学习 Samuel (2024) — 理解双重 Schubert 多项式的乘法公式
\item \textbf{然后}: 学习 Kirillov \& Maeno (1996) — 理解量子版本的基本定义
\item \textbf{最后}: 学习 Mihalcea — 理解量子上同调中 positivity 的推广
\end{enumerate}

\textbf{前置基础知识}:
\begin{itemize}
\item 对称群 $S_n$ 和排列的基本理论
\item Schubert 多项式的定义和性质
\item 量子上同调的基本概念
\item 除差算子 (divided difference operators)
\item 等变量子上同调
\end{itemize}

\section{最终目标}

尝试证明 Gao \& Xiong 结果的量子版本，即在等变量量子上同调的框架下，证明相应的 positivity 性质。

\userannotation{学习路径}{先学 GaoXiong，再学 Samuel，然后 KirillovMaeno，最后 Mihalcea}
\userannotation{新增文献}{添加了 AndersonFulton, Li (2023x2), Buch 作为补充参考}
